\documentclass[conference]{IEEEtran}

\usepackage[utf8]{inputenc}
\usepackage{amsmath,amssymb,amsfonts}
\usepackage{graphicx}
\usepackage{url}
\usepackage{booktabs}
\usepackage{enumitem}
\usepackage{hyperref}
\usepackage{cite}

% TARGET VENUE: SECRYPT 2026 (primary), IEEE TrustCom 2026 (secondary)
% CONTRIBUTION TYPE: System + prototype paper
% CLAIM SCOPE: Prototype architecture only; no benchmark claims without measurement

\begin{document}

\title{PVE-ZKP: Zero-Knowledge Proof-Backed Input Validation for Browser-Based Forms}

\author{
  \IEEEauthorblockN{[Author 1 Name]}
  \IEEEauthorblockA{[Institution] \\[City, Country] \\[email@domain.com]}
}

\maketitle

\begin{abstract}
Web applications routinely outsource input validation to third-party services that
see and store full user submissions, creating avoidable privacy and trust risks.
This paper presents PVE-ZKP, a prototype architecture in which a browser submits
form data to an external validation service that runs five zero-knowledge
policy checks over the inputs, generates a Groth16 proof of policy satisfaction,
and stores the proof under a short identifier. The browser forwards only this
proof identifier to the receiving system, which fetches the proof, verifies it
with \texttt{snarkjs}, and accepts or rejects the submission based solely on
proof validity, without access to raw form inputs. We implement PVE-ZKP for a
healthcare-style registration form using Circom 2.1, snarkjs~0.7, and Node.js,
with policy circuits for format/structure, semantic ranges, set membership,
hidden-field integrity, and path-boundary checks. Our prototype demonstrates
that zero-knowledge input validation can be integrated into familiar
browser--service--backend workflows while preserving data minimization and
trust separation. We discuss security properties, limitations, and how PVE-ZKP
relates to adjacent work on ZK circuit analysis and secure aggregation.
\end{abstract}

\section{Introduction}

Web applications collect sensitive user data---including identifiers, demographic
attributes, and medical information---through HTML forms. Today, validation of
these inputs is typically handled by application backends or third-party
services that receive raw submissions, evaluate policy rules, and return an
accept/reject decision. This design requires both the validation service and the
receiving system to see and store full inputs, enlarging the attack surface and
placing strong trust burdens on intermediaries.

Zero-knowledge proof (ZKP) systems allow a prover to convince a verifier of the
truth of a statement without revealing the underlying witness~\cite{gmr85}.
Recent work has applied ZKPs to Web3 rollups, private transactions, and circuit
security analysis~\cite{chaliasos2024sok,zkap,zkfuzz}, but comparatively little
attention has been paid to using ZKPs for \emph{input validation} in classical
web workflows such as browser form submission.

This paper explores the following question: \emph{Can a browser outsource input
validation to an untrusted service in such a way that the final receiving system
can verify correct validation without learning raw inputs?} To answer this, we
prototype a concrete system called PVE-ZKP (Proof-of-Validated-Execution with
zero-knowledge proofs) and evaluate its suitability as a building block for
privacy-preserving input pipelines.

The contributions of this paper are:
\begin{enumerate}[leftmargin=*]
  \item \textbf{Architecture.} We define a browser--validator--receiver
    architecture in which a validation service at port~3001 runs five Circom
    policy circuits over form inputs, generates Groth16 proofs, and stores them
    under opaque identifiers. The browser sends only a proof identifier onward
    to the receiving system at port~3002, which fetches and verifies the proof
    with \texttt{snarkjs} before accepting or rejecting the submission.
  \item \textbf{Policy circuits.} We design and implement five reusable policy
    circuits that cover format/structure, semantic range, set membership,
    hidden-field integrity via commitments, and path-boundary checks, derived
    from common validation requirements in healthcare-style forms.
  \item \textbf{Prototype implementation.} We implement PVE-ZKP using Circom
    2.1, snarkjs~0.7, circomlib~2.0.5, and Node.js~18, providing build scripts
    for circuit compilation, trusted setup, proof generation, and verification,
    as well as an end-to-end test harness.
  \item \textbf{Security discussion and limitations.} We analyze the trust and
    privacy properties of the architecture, highlight limitations such as
    single-party trusted setup and lack of performance benchmarking
    [VERIFY], and discuss how PVE-ZKP complements circuit-analysis tools such as
    ZKAP and zkFuzz.
\end{enumerate}

The remainder of this paper is organized as follows. Section~II describes the
system model and threat assumptions. Section~III presents the PVE-ZKP
architecture and message flow. Section~IV details the five policy circuits.
Section~V describes the implementation. Section~VI discusses security properties
and limitations. Section~VII relates PVE-ZKP to adjacent work, and
Section~VIII concludes.

\section{System Model and Threat Assumptions}

PVE-ZKP targets web applications in which a user fills out a browser-based form
and a backend system must decide whether to accept or reject the submission
according to validation policies. We consider three logical components:

\begin{itemize}[leftmargin=*]
  \item \textbf{Browser client.} Runs in the user's environment and submits
    form data over HTTPS to a validation service.
  \item \textbf{Validation service.} A potentially third-party service
    (\texttt{validator.js}) listening on port~3001 that receives raw form data,
    runs ZK policy checks, and stores resultant proofs.
  \item \textbf{Receiving system.} An application backend
    (\texttt{receiving\_system.js}) listening on port~3002 that ultimately
    accepts or rejects submissions after verifying proofs.
\end{itemize}

We assume standard network adversaries who may eavesdrop or modify traffic
outside of TLS, and a malicious or curious validation service that attempts to
learn as much as possible about user inputs. The receiving system follows its
verification logic but may be operated by a different organization than the
validator. Our goals are:

\begin{itemize}[leftmargin=*]
  \item \textbf{G1 (Data minimization).} The receiving system should not need to
    see raw form inputs in order to decide whether to accept a submission.
  \item \textbf{G2 (Trust separation).} The system should not have to fully
    trust the validation service; instead, it should verify ZK proofs.
  \item \textbf{G3 (Policy expressiveness).} Common input validation patterns
    should be expressible as Circom circuits without undue complexity.
  \item \textbf{G4 (Deployability).} The architecture should integrate with
    existing browser and backend technologies.
\end{itemize}

\section{PVE-ZKP Architecture}

Figure~\ref{fig:architecture} conceptually illustrates the architecture; we
summarize the interaction in text here. The core workflow is:

\begin{enumerate}[leftmargin=*]
  \item \textbf{Browser \textrightarrow{} validation service.} The browser posts
    JSON-encoded form data to the validation service at port~3001.
  \item \textbf{ZK policy checks.} The validation service maps form fields to
    inputs of five Circom circuits representing format, range, membership,
    integrity, and boundary policies, and runs Groth16 proof generation for each
    applicable policy.
  \item \textbf{Proof storage and identifier.} The validation service stores the
    resulting proofs in an in-memory store keyed by a short, random identifier
    \texttt{proofId}, and returns this identifier to the browser.
  \item \textbf{Browser \textrightarrow{} receiving system.} The browser submits
    an application request that includes only \texttt{proofId} (and any
    non-sensitive metadata) to the receiving system at port~3002. Raw form
    values are \emph{not} forwarded.
  \item \textbf{Proof fetch and verification.} The receiving system calls the
    validation service's proof API, e.g., \texttt{GET /proof/:id}, to fetch the
    stored proof(s), and verifies them using the pre-generated verification keys
    and \texttt{snarkjs}.
  \item \textbf{Accept / Reject.} If all required proofs verify, the receiving
    system accepts and processes the submission; otherwise, it rejects and
    returns an error.
\end{enumerate}

This workflow is summarized textually as:
\begin{quote}
Browser (client) $\rightarrow$ \texttt{validator.js}:3001 (proof generation,
proof storage under \texttt{proofId}) $\rightarrow$ browser (forward
\texttt{proofId}) $\rightarrow$ \texttt{receiving\_system.js}:3002 (fetch and
verify proofs) $\rightarrow$ accept / reject.
\end{quote}

In the prototype, proof storage is an in-memory map, which is suitable for
experiments but would need to be replaced with a durable store such as Redis
with appropriate access control and time-to-live in production.

\section{Policy Circuits}

PVE-ZKP implements five Circom circuits aligned with the ``Policy Circuits''
section of the open-source prototype:

\begin{itemize}[leftmargin=*]
  \item \textbf{\texttt{format\_check}} (Format/Structure): ensures that a
    fixed-length character array (e.g., a 7-character postcode) consists of
    allowed ASCII characters and obeys simple structure rules.
  \item \textbf{\texttt{threshold\_check}} (Semantic Range): enforces that a
    numeric field such as age falls within a minimum or maximum range.
  \item \textbf{\texttt{membership\_check}} (Approved Set): checks that a
    value belongs to an approved list, e.g., country codes or insurance types.
  \item \textbf{\texttt{integrity\_check}} (Hidden Field Integrity): binds a
    hidden token, a server-side secret, and a session identifier into a Poseidon
    hash commitment that the receiver can later recompute to ensure the form has
    not been tampered with.
  \item \textbf{\texttt{boundary\_check}} (Path Boundary): ensures that a
    path-like string (e.g., a referral URL prefix) stays within an allowed
    boundary.
\end{itemize}

These circuits are written in Circom 2.1 and rely on circomlib 2.0.5. The
implementation follows the prototype's README, which summarizes private and
public inputs for each policy.

Due to space constraints, full circuit listings are omitted; the open-source
repository contains the complete Circom files along with build scripts.

\section{Implementation}

The prototype implementation consists of:
\begin{itemize}[leftmargin=*]
  \item Circom 2.1 circuit files for the five policies.
  \item A \texttt{build\_circuits.sh} script that downloads a powers-of-tau file,
    compiles circuits, runs Groth16 trusted setup, and exports verification keys.
  \item A Node.js validation service (\texttt{validator.js}) on port~3001 that
    exposes HTTP endpoints for proof generation and maintains an in-memory proof
    store keyed by \texttt{proofId}.
  \item A Node.js receiving system (\texttt{receiving\_system.js}) on port~3002
    that fetches proofs by identifier and verifies them using \texttt{snarkjs}.
  \item A browser-based demo form (HTML + JavaScript) that submits form data to
    the validator and then forwards \texttt{proofId} to the receiver.
\end{itemize}

The prototype is packaged with a quick-start workflow: install dependencies,
build circuits, run end-to-end tests over sample inputs, and start both
services. This mirrors common Node.js deployment practices, making it easier for
non-cryptography specialists to experiment with ZK-backed validation.

\section{Security Properties and Limitations}

PVE-ZKP aims to provide data minimization (G1) and trust separation (G2) with
respect to input validation, assuming the underlying cryptographic primitives
and Circom tooling are correct.

\subsection{Security Properties}

\textbf{Correctness.} For honest clients and a non-malicious validator, valid
inputs that satisfy all policy circuits produce proofs that the receiving system
accepts. This follows from Groth16 completeness and correct circuit
implementation.

\textbf{Soundness against malicious clients.} A malicious browser that attempts
to submit invalid inputs must convince the validator to generate proofs for
unsatisfied policies. Assuming circuits are sound and the Groth16 proof system
is instantiated correctly, this should be infeasible except with negligible
probability.

\textbf{Data minimization.} The receiving system bases its decision solely on
verified proofs and does not receive raw form values from the browser. It may
still learn some information through the fact that a proof exists (e.g., that
inputs satisfy policy ranges), which is inherent to the validation outcome.

\textbf{Trust separation.} The receiving system does not have to fully trust the
validation service: whether or not the validator behaves honestly, the receiver
can independently verify proofs before accepting submissions.

\subsection{Limitations}

\textbf{Trusted setup.} The prototype uses a single-party Groth16 trusted setup
script. A production deployment would require a multi-party ceremony to avoid
trusted setup concerns.

\textbf{Proof storage.} Proofs are stored in an in-memory map in the validation
service. This is not durable and is unsuitable for high-availability scenarios;
moreover, access control and retention policies are not enforced.

\textbf{Performance and cost.} We do not yet report empirical measurements for
proving and verification time, memory usage, or proof sizes; these must be
measured for realistic deployments [VERIFY].

\textbf{Side-channel leakage.} Timing, error messages, and traffic patterns may
leak additional information about inputs beyond the ZK proofs themselves. The
prototype does not implement mitigations such as constant-time paths.

\textbf{Policy coverage.} The five policy circuits cover common validation
patterns but do not constitute a complete validation language; extending PVE-ZKP
to richer policies (e.g., cross-field dependencies) is left as future work.

\section{Related Work}

Chaliasos et al.~\cite{chaliasos2024sok} survey security vulnerabilities across
SNARK-based systems, but do not focus on browser form validation. ZKAP~\cite{zkap}
performs static analysis of ZK circuits to find underconstrained signals and
unsafe compositions, while zkFuzz~\cite{zkfuzz} introduces a fuzzing framework
for ZK circuits. These tools are complementary to PVE-ZKP: they help ensure the
soundness of the policy circuits used in our architecture.

Adjacent lines of work consider input validation for aggregated statistics and
federated learning. VPAS and LZKSA-style schemes~\cite{lzksa,risefl} use ZKPs or
post-quantum secure proofs to ensure that individual client contributions to an
aggregate model satisfy policy constraints without revealing the contributions
themselves. In contrast, PVE-ZKP targets single-form submissions in classical
web applications.

Finally, secure web architectures such as privacy-preserving validation servers
have explored non-ZK techniques (e.g., differential privacy) to reduce data
exposure; PVE-ZKP can be viewed as a cryptographic strengthening of these ideas
by providing explicit proofs of policy satisfaction.

\section{Conclusion}

PVE-ZKP shows that zero-knowledge proofs can be integrated into familiar
browser-based input workflows to reduce trust in validation services and minimize
exposure of raw user data to backend systems. By separating the roles of
validator and receiver and basing acceptance decisions on verifiable ZK policy
proofs, the architecture offers a concrete path toward privacy-preserving input
validation in domains such as healthcare forms. Future work includes measuring
performance overheads, extending the policy circuit library, integrating
circuit-analysis tools into the development workflow, and exploring
multi-tenant deployments with stronger setup ceremonies.

\bibliographystyle{IEEEtran}

\begin{thebibliography}{99}

\bibitem{gmr85}
S.~Goldwasser, S.~Micali, and C.~Rackoff,
``The knowledge complexity of interactive proof systems,''
\textit{SIAM Journal on Computing}, vol.~18, no.~1, pp.~186--208, 1989.

\bibitem{chaliasos2024sok}
S.~Chaliasos \textit{et al.},
``SoK: What don't we know? Understanding security vulnerabilities in SNARKs,''
in \textit{Proc. 33rd USENIX Security Symposium}, 2024.

\bibitem{zkap}
H.~Wen \textit{et al.},
``Practical security analysis of zero-knowledge proof circuits,''
in \textit{Proc. 33rd USENIX Security Symposium}, 2024.

\bibitem{zkfuzz}
J.~Yang \textit{et al.},
``zkFuzz: Foundation and framework for effective fuzzing of zero-knowledge circuits,''
arXiv:2504.11961, 2025.

\bibitem{lzksa}
Z.~Lu \textit{et al.},
``LZKSA: Lattice-based special zero-knowledge proofs for secure aggregation,''
in \textit{Proc. ACM CCS}, 2025.

\bibitem{risefl}
X.~Zhu \textit{et al.},
``Secure and verifiable data collaboration with low-cost zero-knowledge proofs,''
arXiv:2311.15310, 2023.

\end{thebibliography}

\end{document}
